% Isolierte Fixture fuer die Regeln E-I aus Tabelle 1.
%
% Die Datei verwendet die produktiven Formatdefinitionen mit kontrolliertem
% Inhalt, ohne main.tex zu veraendern.
\documentclass[
  12pt,
  a4paper,
  oneside,
  headings=normal,
  numbers=noendperiod,
  bibliography=totoc
]{scrreprt}

\newcommand{\MainLanguage}{ngerman}
\newcommand{\ThesisType}{Compliance-Test}
\newcommand{\ThesisTitle}{Tabelle-1-Compliance}
\newcommand{\AuthorName}{Compliance-Test}
\newcommand{\MatriculationNumber}{0000000}
\newcommand{\DegreeProgramme}{Teststudiengang}
\newcommand{\Department}{Testfachbereich}
\newcommand{\FirstExaminer}{Erstprüfung}
\newcommand{\SecondExaminer}{Zweitprüfung}
\newcommand{\SubmissionDate}{1. Januar 2026}

\usepackage{hfh-formatierungen}
\setcounter{secnumdepth}{3}
\setcounter{tocdepth}{3}
\addbibresource{fixtures/compliance-references.bib}

\begin{document}

\chapter{Überschriftenabstände}

Dieser Absatz prüft den Abstand von einer Hauptüberschrift zum unmittelbar
folgenden Fließtext. Tabelle 1 verlangt hierfür 18 Punkt nach der
Überschrift.

\section{Zweite Ebene nach Fließtext}

Dieser Absatz prüft den Abstand vor und nach der zweiten Gliederungsebene.
Tabelle 1 verlangt 18 Punkt davor und 12 Punkt danach.

\subsection{Dritte Ebene nach Fließtext}

Dieser Absatz prüft dieselben Übergänge für die dritte Ebene. Schriftgröße,
Fettdruck und einzeiliger Satz entsprechen den Tabelle-1-Regeln.

\subsubsection{Vierte Ebene nach Fließtext}

Dieser Absatz prüft die vierte Ebene, die bei umfangreichen Arbeiten optional
aktiviert werden kann.

\section{Direkter Übergang zwischen Ebenen}
\subsection{Zweite auf dritte Ebene}
\subsubsection{Dritte auf vierte Ebene}

Aufeinanderfolgende Überschriften dürfen nicht durch den normalen
Absatzabstand verfälscht werden. Diese Folge macht den Übergang sichtbar.

\chapter{Fußnoten innerhalb eines Kapitels}

Die erste Markierung verweist auf eine mehrzeilige Fußnote.\footnote{Diese
erste Fußnote beginnt mit einem Großbuchstaben, ist bewusst mehrzeilig und
endet mit dem vorgeschriebenen Punkt.} Mehrere Noten auf derselben Seite
machen sowohl den einzeiligen Innensatz als auch den Abstand zwischen den
Noten sichtbar.\footnote{Diese zweite Fußnote enthält zusätzlichen Text, damit
mindestens zwei Zeilen entstehen und der Grundlinienabstand zuverlässig
beurteilt werden kann.}

Ein weiterer Absatz erzeugt eine dritte Fußnote.\footnote{Diese dritte
Fußnote prüft erneut den ungefähr vier Punkt großen Zusatzabstand zur
vorhergehenden Fußnote.} Der Text danach stellt sicher, dass der Fußnotenblock
nicht mit dem Haupttext kollidiert.

\chapter{Fußnoten nach einem Kapitelwechsel}

Die Nummerierung darf beim Kapitelwechsel nicht neu beginnen. Deshalb muss
diese Fußnote die Nummer vier tragen.\footnote{Diese vierte Fußnote belegt die
fortlaufende Nummerierung über die Kapitelgrenze hinweg.} Eine zweite Note in
diesem Kapitel muss folgerichtig die Nummer fünf erhalten.\footnote{Diese
fünfte Fußnote schließt den positiven Fußnotentest mit einem Punkt ab.}

\chapter{Tabellen und Abbildungen}

Tabelle~\ref{tab:overview} und Abbildung~\ref{fig:process} werden vor ihrem
Auftreten ausdrücklich im Fließtext referenziert.

\begin{hfhtable}
  \caption{Mehrzeilige Übersicht der geprüften Tabellenmerkmale}
  \label{tab:overview}
  \begin{tabularx}{\textwidth}{@{}Yp{2.4cm}p{2.4cm}@{}}
    \toprule
    \textbf{Prüfmerkmal} & \textbf{Sollwert} & \textbf{Status} \\
    \midrule
    Schriftgröße und Zeilenabstand innerhalb eines längeren Tabellenfeldes &
    10 Punkt & einzeilig \\
    Abstand zwischen Beschriftung und Tabellenkörper &
    6 Punkt & sichtbar \\
    \bottomrule
  \end{tabularx}
  \source{Eigene Darstellung.}
\end{hfhtable}

\begin{hfhfigure}
  \fbox{%
    \begin{minipage}[c][3.2cm][c]{0.68\textwidth}
      \centering\footnotesize\setstretch{1}
      Eingabe\\Prüfung\\Freigabe
    \end{minipage}}
  \caption{Einzeilig gesetzter Prozess innerhalb einer Abbildung}
  \label{fig:process}
  \source{Eigene Darstellung.}
\end{hfhfigure}

Der Text nach der Abbildung prüft den Abstand von ungefähr 12 Punkt nach der
Beschriftungs- und Quellenzone.

\chapter{Globale Nummerierung}

Tabelle~\ref{tab:second} und Abbildung~\ref{fig:second} müssen trotz des neuen
Kapitels global als Nummer zwei erscheinen.

\begin{hfhtable}
  \caption{Zweite Tabelle in einem neuen Kapitel}
  \label{tab:second}
  \begin{tabularx}{\textwidth}{@{}Yp{3cm}@{}}
    \toprule
    \textbf{Merkmal} & \textbf{Ergebnis} \\
    \midrule
    Kapitelunabhängiger Tabellenzähler & Tabelle 2 \\
    \bottomrule
  \end{tabularx}
  \source{Eigene Darstellung.}
\end{hfhtable}

\begin{hfhfigure}
  \fbox{\rule{0pt}{2.2cm}\rule{0.58\textwidth}{0pt}}
  \caption{Zweite Abbildung in einem neuen Kapitel}
  \label{fig:second}
  \source{Eigene Darstellung.}
\end{hfhfigure}

\chapter{Quellenverzeichnisse}

Die Literatur wurde absichtlich in nichtalphabetischer Reihenfolge zitiert
\parencite{zimmer2022,adler2020,becker2021}. Der Bibliografiestil muss die
Einträge
trotzdem alphabetisch nach Autor ausgeben.

\printbibliography[heading=bibintoc,title={Quellenverzeichnis}]

\addchap{Rechtsprechungsverzeichnis}

\begin{hfhsourceentries}
  \hfhsourceentry{Arbeitsgericht Berlin, Urteil vom 12. Januar 2020,
    Aktenzeichen 1 Ca 100/20, veröffentlichte Sammlung 2020, S. 10--18.}
  \hfhsourceentry{Bundesgerichtshof, Beschluss vom 3. März 2021,
    Aktenzeichen II ZB 2/21, amtliche Sammlung 2021, S. 100--112.}
  \hfhsourceentry{Verwaltungsgericht Hamburg, Urteil vom 9. September 2022,
    Aktenzeichen 3 K 300/22, Fachzeitschrift 2023, S. 30--42.}
\end{hfhsourceentries}

\phantomsection
\label{LastNumberedPage}

\end{document}
